\documentclass[12pt]{article}
\usepackage[a4paper,margin=1in]{geometry}
\usepackage[utf8]{inputenc}
\usepackage[T1]{fontenc}
\usepackage{amsmath,amssymb}
\usepackage{textcomp}
\usepackage{parskip}
\begin{document}
\begin{center}
{\Large \textbf{2000 STEP III}}
\end{center}

\subsection*{Page 1}
\begin{flushleft}
STEP III, 2000 2\par
\bigskip\textbf{Section A: Pure Mathematics}\par
\bigskip\textbf{Question 1.} Sketch on the same axes the two curves C and C , given by\par
1 2\par
\[C : xy = 1 ,\]
\[C : x^2 y^2 = 2 .\]
-\par
The curves intersect at P and Q. Given that the coordinates of P are (a, b) (which you need\par
not evaluate), write down the coordinates of Q in terms of a and b.\par
The tangent to C through P meets the tangent to C through Q at the point M , and the\par
1 2\par
tangent to C through P meets the tangent to C through Q at N . Show that the coordinates\par
2 1\par
of M are ( b, a) and write down the coordinates of N .\par
-\par
Show that P M QN is a square.\par
\bigskip\textbf{Question 2.} Use the substitution x = 2 cos θ to evaluate the integral\par
-\par
2 x 1 2\par
- dx.\par
Z µ 3 x ¶\par
\[3/2\]
-\par
\[Show that, for a < b,\]
\[q x a 2 (b a)(\pi + 3\sqrt\{3\} 6)\]
\[- dx = - - ,\]
Z µ b x ¶ 12\par
p\par
-\par
\[where p = (3a + b)/4 and q = (a + b)/2.\]
\[\pi\]
\bigskip\textbf{Question 3.} Given that α = ei 3 , prove that 1 + α2 = α.\par
A triangle in the Argand plane has vertices A, B, and C represented by the complex numbers p,\par
qα2 and rα respectively, where p, q and r are positive real numbers. Sketch the triangle ABC.\par
-\par
Three equilateral triangles ABL, BCM and CAN (each lettered clockwise) are erected on\par
sides AB, BC and CA respectively. Show that the complex number representing N is\par
(1 α)p α2r and find similar expressions for the complex numbers representing L and M .\par
- -\par
Show that lines LC, M A and N B all meet at the origin, and that these three line segments\par
have the common length p + q + r.\par
\end{flushleft}
\subsection*{Page 2}
\begin{flushleft}
STEP III, 2000 3\par
\bigskip\textbf{Question 4.} The function f (x) is defined by\par
x(x 2)(x a)\par
\[f (x) = - - .\]
\[x^2 1\]
-\par
Prove algebraically that the line y = x + c intersects the curve y = f (x) if a > 1, but there\par
| |\par
\[are values of c for which there are no points of intersection if a < 1.\]
| |\par
\[Find the equation of the oblique asymptote of the curve y = f (x). Sketch the graph in the\]
\[two cases (i) a < 1 ; and (ii) 1 < a < 1 . (You need not calculate the turning points.)\]
- - - 2\par
a b\par
\bigskip\textbf{Question 5.} Given two non-zero vectors a = 1 and b = 1 we define ∆(a, b) by ∆(a, b) = a b a b .\par
µa 2¶ µb 2¶ 1 2 - 2 1\par
Let A, B and C be points with position vectors a, b and c, respectively, no two of which are\par
parallel. Let P , Q and R be points with position vectors p, q and r, respectively, none of\par
which are parallel.\par
\textbf{(i)} Show that there exists a 2 2 matrix M such that P and Q are the images of A and\par
\[\times\]
B under the transformation represented by M.\par
\textbf{(ii)} Show that ∆(a, b) c + ∆(c, a) b + ∆(b, c) a = 0.\par
Hence, or otherwise, prove that a necessary and sufficient condition for the points P ,\par
Q, and R to be the images of points A, B and C under the transformation represented\par
by some 2 2 matrix M is that\par
\[\times\]
\[∆(a, b) : ∆(b, c) : ∆(c, a) = ∆(p, q) : ∆(q, r) : ∆(r, p) .\]
\bigskip\textbf{Question 6.} Given that\par
\[x^4 + px2 + qx + r = (x^2 ax + b)(x^2 + ax + c),\]
-\par
express p, q and r in terms of a, b and c.\par
\[Show also that a^2 is a root of the cubic equation\]
\[u^3 + 2pu2 + (p^2 4r)u q^2 = 0.\]
- -\par
Explain why this equation always has a non-negative root, and verify that u = 9 is a root in\par
\[the case p = 1, q = 6, r = 15 .\]
- -\par
Hence, or otherwise, express\par
\[y^4 8y3 + 23y2 34y + 39\]
- -\par
as a product of two quadratic factors.\par
\end{flushleft}
\subsection*{Page 3}
\begin{flushleft}
STEP III, 2000 4\par
\bigskip\textbf{Question 7.} Given that\par
1 1 1 1\par
\[e = 1 + + + + + + ,\]
1! 2! 3! \cdots r! \cdots\par
use the binomial theorem to show that\par
n\par
\[1 + < e\]
µ n ¶\par
for any positive integer n.\par
The product P(n) is defined, for any positive integer n, by\par
3 5 9 2n + 1\par
\[P(n) = . . . .\]
2 · 4 · 8 · · 2n\par
Use the arithmetic-geometric mean inequality,\par
\[a 1 + a 2 + \cdots + a n > (a a . . . a ) n 1 ,\]
1 2 n\par
n \cdots\par
\[to show that P(n) < e for all n .\]
Explain briefly why P(n) tends to a limit as n . Show that this limit, L, satisfies\par
\to \infty\par
\[2 < L 6 e.\]
\bigskip\textbf{Question 8.} The sequence a is defined by a = 1 , a = 1 , and\par
n 0 1\par
\[1 + a^2\]
\[a =\]
n-1\par
\[(n > 2).\]
n\par
a n-2\par
Prove by induction that\par
\[a n = 3a n-1 a n-2 (n > 2).\]
-\par
Hence show that\par
α2n-1\par
+ α\par
-(2n-1)\par
\[a = (n > 1),\]
n\par
\[\sqrt\{5\}\]
\[1 + \sqrt\{5\}\]
\[where α = .\]
\end{flushleft}
\subsection*{Page 4}
\begin{flushleft}
STEP III, 2000 5\par
\bigskip\textbf{Section B: Mechanics}\par
\bigskip\textbf{Question 9.} Two small discs of masses m and µm lie on a smooth horizontal surface. The disc of mass µm\par
is at rest, and the disc of mass m is projected towards it with velocity u. After the collision,\par
the disc of mass µm moves in the direction given by unit vector n. The collision is perfectly\par
elastic.\par
2u.n\par
\textbf{(i)} Show that the speed of the disc of mass µm after the collision is .\par
1 + µ\par
\textbf{(ii)} Given that the two discs have equal kinetic energy after the collision, find an expres-\par
\[sion for the cosine of the angle between n and u and show that 3 \sqrt\{8\} 6 µ 6 3 + \sqrt\{8\}.\]
-\par
\bigskip\textbf{Question 10.} A sphere of radius a and weight W rests on horizontal ground. A thin uniform beam of weight\par
3\sqrt3 W and length 2a is freely hinged to the ground at X, which is a distance \sqrt3 a from the\par
point of contact of the sphere with the ground. The beam rests on the sphere, lying in the\par
same vertical plane as the centre of the sphere. The coefficients of friction between the beam\par
and the sphere and between the sphere and the ground are µ and µ respectively.\par
1 2\par
Given that the sphere is on the point of slipping at its contacts with both the ground and the\par
beam, find the values of µ and µ .\par
1 2\par
\bigskip\textbf{Question 11.} A thin beam is fixed at a height 2a above a horizontal plane. A uniform straight rod ACB\par
of length 9a and mass m is supported by the beam at C. Initially, the rod is held so that it\par
is horizontal and perpendicular to the beam. The distance AC is 3a, and the coefficient of\par
friction between the beam and the rod is µ.\par
The rod is now released. Find the minimum value of µ for which B strikes the horizontal\par
plane before slipping takes place at C.\par
\end{flushleft}
\subsection*{Page 5}
\begin{flushleft}
STEP III, 2000 6\par
\bigskip\textbf{Section C: Statistics}\par
\bigskip\textbf{Question 12.} In a lottery, any one of N numbers, where N is large, is chosen at random and independently\par
for each player by machine. Each week there are 2N players and one winning number is drawn.\par
Write down an exact expression for the probability that there are three or fewer winners in a\par
week, given that you hold a winning ticket that week. Using the fact that\par
n\par
a\par
1 e\par
-a\par
µ - n ¶ \approx\par
for n much larger than a, or otherwise, show that this probability is approximately 2 .\par
Discuss briefly whether this probability would increase or decrease if the numbers were chosen\par
by the players.\par
Show that the expected number of winners in a week, given that you hold a winning ticket\par
that week, is 3 N\par
-1.\par
-\par
\bigskip\textbf{Question 13.} A set of n dice is rolled repeatedly. For each die the probability of showing a six is p. Show\par
that the probability that the first of the dice to show a six does so on the rth roll is\par
qnr(q -n 1)\par
-\par
\[where q = 1 p.\]
-\par
Determine, and simplify, an expression for the probability generating function for this distri-\par
bution, in terms of q and n. The first of the dice to show a six does so on the Rth roll. Find\par
\[the expected value of R and show that, in the case n = 2, p = 1/6, this value is 36/11.\]
Show that the probability that the last of the dice to show a six does so on the rth roll is\par
(1 qr) n (1 qr-1) n .\par
- - -\par
Find, for the case n = 2, the probability generating function. The last of the dice to show a\par
\[six does so on the Sth roll. Find the expected value of S and evaluate this when p = 1/6.\]
\end{flushleft}
\subsection*{Page 6}
\begin{flushleft}
STEP III, 2000 7\par
\bigskip\textbf{Question 14.} The random variable X takes only the values x and x (where x = x ), and the random\par
1 2 1 2\par
\[variable Y takes only the values y and y (where y = y ). Their joint distribution is given\]
1 2 1 2\par
by\par
\[P(X = x , Y = y ) = a ; P(X = x , Y = y ) = q a ; P(X = x , Y = y ) = p a .\]
1 1 1 2 2 1\par
- -\par
\[Show that if E(XY ) = E(X)E(Y ) then\]
\[(a pq)(x x )(y y ) = 0.\]
1 2 1 2\par
- - -\par
Hence show that two random variables each taking only two distinct values are independent\par
\[if E(XY ) = E(X)E(Y ).\]
Give a joint distribution for two random variables A and B, each taking the three values 1,\par
-\par
\[0 and 1 with probability 1 , which have E(AB) = E(A)E(B), but which are not independent.\]
\end{flushleft}
\end{document}
